\chapter{Introduction, Aims and Objectives}
\label{ch:introduction}

\section{Introduction to Problem}

For more than fifty years, the password has been the main way people prove
who they are on the web. It has not aged well. People are asked to make
up, remember, and regularly change dozens of secret words for accounts they
barely use, so they do what any person under that kind of pressure would
do: they pick easy passwords, and they use the same one everywhere. The
2024 Verizon Data Breach Investigations Report looked at 10,626 confirmed
data breaches and found that stealing or guessing login details was the
single most common way attackers got in, the starting point in 38\% of
cases, with phishing close behind at 15\%, and human error playing a part
in 68\% of all breaches \cite{verizon2024dbir}. The problem is not that
people are careless. The problem is that passwords ask ordinary users to
be perfect guardians of a secret that is easy to steal by fake emails, easy
to reuse, and expensive to change, while proving nothing about whether the
person typing it is really who they say they are.

The tech industry has answered with a wave of "passwordless" login methods:
one-time codes sent by text or email, magic links, QR-code login, and,
more and more, passkeys built on a standard called FIDO2/WebAuthn. Each of
these fixes part of the problem but leaves another part open. A text
message code can still be phished, and it depends on the phone network
being secure. A magic link removes the typing but adds a second thing the
user must switch to, their email inbox, and now the login is only as safe
as that inbox. QR-code login means aiming a camera at a screen, which is
awkward. Passkeys get the hard cryptography right, tying a secret key to
one specific website, but the user still has to notice a pop-up on their
screen, pick the right device, and deal with quirks that differ from one
browser or phone to the next. None of these widely used methods let a
person log in just by having their trusted phone nearby and giving a quick
nod of approval.

This project tries to close that gap. It asks whether a sound that people
cannot hear, sent between a laptop and a phone, combined with strong
cryptography, can give a login experience with no typing beyond a
username, no aiming a camera, and no switching to a second app, while
still standing up to the same attacks that make passwords and their
replacements weak: phishing, replaying a stolen login attempt, and
stealing login details from far away.

\section{Introduction to Project, Aim and Objectives}

\subsection*{Project Aim}

The aim of this project is to design, build, and test Echo, a passwordless
login system for websites. A laptop and a phone that the user has already
registered talk to each other using a sound too high-pitched for most
people to hear, backed by public-key cryptography. The project also checks,
through real testing, whether this idea is both cryptographically sound and
reliable enough for everyday use.

\subsection*{Objectives}

\begin{enumerate}[label=\arabic*.]
  \item Look at why passwords fail in practice, and compare the trade-offs
  made by the passwordless methods already in use (one-time codes, magic
  links, QR-code login, and FIDO2/WebAuthn passkeys), to explain why Echo
  is designed the way it is.
  \item Build a sound-based communication layer, using a library called
  \texttt{ggwave}, that reliably sends a one-time code from a laptop's
  speakers to a phone's microphone at 18 to 20\,kHz (too high for most
  people to hear), with a normal audible backup for hardware that cannot
  play those high frequencies.
  \item Build the Echo server: registering users and phones, creating and
  managing one-time codes, checking digital signatures, sending live
  updates over WebSocket, and letting users recover their account, all
  backed by a SQLite database.
  \item Build the phone app (a Progressive Web App, meaning it installs
  from a browser): it listens through the microphone, decodes the sound,
  shows an approve or deny screen, creates a private signing key that
  never leaves the phone, and signs login requests.
  \item Build a set of automated tests that check, without needing real
  audio hardware, that a replayed login sound, a forged signature, an
  expired code, and someone else's phone are all correctly rejected by the
  server.
  \item Test how well the sound actually works in the real world, at
  different distances (0.3\,m, 1\,m, 2\,m) and in different amounts of
  background noise (quiet, speech, music), aiming for at least 90\%
  success at 1\,m in a quiet room.
  \item Build and test a way for users to recover their account so they
  are never locked out for good, along with basic safety features like
  limiting repeated attempts.
\end{enumerate}

\section{Research Questions}

The testing in this report is built around four questions:

\begin{enumerate}[label=RQ\arabic*.]
  \item Can an everyday laptop speaker and an everyday phone microphone
  reliably send a short piece of data at a near-silent pitch (18 to 20\,kHz)
  across normal login distances (0.3 to 2\,m) and in real background noise?
  \item Does tying that sound to a one-time code and a phone-specific
  digital signature stop the classic attacks that break weaker
  passwordless methods, namely replaying a recorded login and faking a
  signature?
  \item How long does the whole process take, from the server creating the
  code to the laptop being logged in, and does that compare well with
  other passwordless methods?
  \item What risks does a sound-based approach still leave open, and how
  should those risks be explained honestly in a project like this, rather
  than glossed over?
\end{enumerate}

\section{Scope}

This project is meant to produce a working system that can be properly
tested, not a finished commercial product. Echo includes: signing up with
a username and registering one phone; logging in using a one-time code
sent as sound; digital signatures and server-side checking; live updates
over WebSocket and cookie-based sessions; a rate-limited way to recover an
account by email; and the ability to remove a registered phone. It
deliberately leaves out: a native phone app (the phone side is a
browser-installed app instead, to save time during the semester); strong
defences against an attacker relaying the sound over a network rather than
being physically present (this risk is discussed, but not solved, since
solving it needs hardware control a web browser does not give);
production-level account features such as email verification or an admin
panel; and letting one account register more than one phone. These
choices are explained fully in Section~\ref{sec:further-work} as things to
build next, not as gaps the report tries to hide.

\section{Project Justification}

Passwordless login is not a small or unusual topic. Major tech companies,
browser makers, and an industry group called the FIDO Alliance have spent
the last ten years pushing toward exactly this goal, and the rollout of
passkeys by Apple, Google, and Microsoft since 2022 shows where the whole
industry is heading. But the mainstream solutions are built to fit large
companies and specific platforms, which leaves an interesting and useful
question for a security project: what is the smallest possible way to
prove "the user's trusted phone is here, and its owner approved this
login," without depending on one company's infrastructure, and how far can
a small, understandable, self-built version of that idea be pushed within
one semester? Building the sound channel, the one-time code system, the
signing process, and the checking logic from scratch, instead of using an
existing WebAuthn library, forces every decision to be explicit and
checkable: what is secret, what is public, what expires, and what an
attacker can and cannot do with a recording, a stolen signature, or a fake
website. That is what makes this a good Cybersecurity capstone project,
and it is why this report spends as much space on the threat model and the
failed-attack tests as it spends on the working demo.

\section{Organization of Chapters}

The rest of this report has four more chapters. Chapter~\ref{ch:literature}
looks at why passwords fail, compares the passwordless methods already in
use, and reviews the sound-based and cryptographic techniques Echo is
built on. Chapter~\ref{ch:methodology} explains the step-by-step way the
system was built, its overall design, the login process itself, and the
technology choices behind each part. Chapter~\ref{ch:design} shows the
actual design of the server, the laptop page, and the phone app, the
automated tests and their results, and the real-world testing of sound
reliability and login speed across distance and noise. Chapter~\ref{ch:conclusion}
draws conclusions against the aims and objectives set out here, states the
project's limitations honestly, and suggests what to build next.
References and supporting appendices, including the project plan and the
full test results, are at the end of the report.
